\documentclass[10pt]{article}

% -- Page geometry --
\usepackage[letterpaper, margin=0.5in, top=0.55in, bottom=0.5in]{geometry}

% -- Core packages --
\usepackage{amsmath, amssymb}
\usepackage{array, tabularx, multirow, makecell}
\usepackage{graphicx}
\usepackage{tikz}
\usepackage{pgfplots}
\pgfplotsset{compat=1.18}
\usepackage{enumitem}
\usepackage{tcolorbox}
\tcbuselibrary{skins, breakable}
\usepackage{fancyhdr}
\usepackage{xcolor}
\usepackage{parskip}

% -- Colors --
\definecolor{tealheader}{HTML}{0097A7}
\definecolor{lightgray}{HTML}{F2F2F2}
\definecolor{warmred}{HTML}{C0392B}
\definecolor{dayblue}{HTML}{2B5C8A}

% -- Header / Footer --
\pagestyle{fancy}
\fancyhf{}
\renewcommand{\headrulewidth}{0pt}
\fancyhead[L]{\textsf{FULL NAME: \rule{2.5in}{0.4pt}}}
\fancyhead[R]{\textsf{BLOCK: \rule{1in}{0.4pt}}}
\fancyfoot[C]{\small\textsf{\thepage}}

% -- Fonts --
\usepackage{helvet}
\renewcommand{\familydefault}{\sfdefault}

% -- Custom commands --
\newcommand{\writespace}[1]{\vspace{#1}}
\newcommand{\writelines}[1]{%
  \par\vspace{4pt}%
  \foreach \n in {1,...,#1}{\noindent\rule{\linewidth}{0.3pt}\vspace{8pt}\par}%
}
\newcommand{\gridpaper}[2]{%
  \begin{tikzpicture}
    \draw[step=0.25cm, gray!35, very thin] (0,0) grid (#1,#2);
    \draw[step=1cm, gray!55, thin] (0,0) grid (#1,#2);
  \end{tikzpicture}%
}
\newcommand{\blankline}{\rule{1.5in}{0.4pt}}
\newcommand{\dayheading}[3]{%
  \begin{tcolorbox}[enhanced, colback=dayblue, colframe=dayblue,
    arc=0pt, left=8pt, right=8pt, top=5pt, bottom=5pt,
    fontupper=\Large\bfseries\sffamily\color{white}]
    Day #1 \;\textbar\; #2
  \end{tcolorbox}
  \vspace{1pt}
  {\small\itshape #3}%
}

% -- Column types --
\newcolumntype{C}[1]{>{\centering\arraybackslash}m{#1}}

% -- tcolorbox styles --
\newtcolorbox{activitybox}[1][]{%
  enhanced, colback=white, colframe=black!40,
  boxrule=0.5pt, arc=0pt, left=5pt, right=5pt,
  top=3pt, bottom=3pt, #1
}

\newtcolorbox{generalizebox}{%
  enhanced, colback=lightgray, colframe=lightgray,
  boxrule=0pt, arc=0pt, left=6pt, right=6pt,
  top=4pt, bottom=4pt,
  borderline west={3pt}{0pt}{tealheader}
}

\newtcolorbox{cerbox}{%
  enhanced, colback=white, colframe=warmred!70,
  boxrule=0.8pt, arc=0pt, left=6pt, right=6pt,
  top=4pt, bottom=4pt,
  title={\sffamily\bfseries\color{warmred} CER (Exit Ticket)},
  coltitle=warmred, colbacktitle=warmred!8,
  attach boxed title to top left={yshift=-2mm, xshift=4mm},
  boxed title style={arc=0pt, boxrule=0pt}
}

\newtcolorbox{sectionlabel}[1][]{%
  enhanced, colback=tealheader!10, colframe=tealheader,
  boxrule=0pt, arc=0pt, left=6pt, right=6pt,
  top=2pt, bottom=2pt,
  borderline west={3pt}{0pt}{tealheader},
  #1
}

% -- Suppress paragraph indent --
\setlength{\parindent}{0pt}
\setlength{\parskip}{3pt}

\begin{document}

% ===================================================================
%   COVER / OBJECTIVES PAGE
% ===================================================================

\begin{center}
  {\LARGE Algebra 2 Lesson: 3-5 Zeros of Polynomial Functions}

  \vspace{2pt}
  {\normalsize\textit{8-Day Student Packet}}
\end{center}

\vspace{-2pt}

% -- Objectives Table --
\begin{center}
\renewcommand{\arraystretch}{1.35}
\begin{tabular}{|>{\centering\arraybackslash}m{1.15in}|m{5.6in}|}
\hline
\textbf{Math Objective} &
Identify the zeros of a function using factoring or synthetic division. Use zeros to graph a polynomial function. \\
\hline
\textbf{Language Objective} &
Students will write what they think about a problem using ``Notice and Wonder''. \\
\hline
\textbf{Essential Understanding} &
The zeros of a polynomial function can be determined using factoring or synthetic division. The zeros can be used to sketch its graph. \\
\hline
\end{tabular}
\end{center}

\vspace{10pt}

% ===================================================================
%   DAY 1
% ===================================================================

\dayheading{1}{What Is a Zero?}{Complete on your laptop, in a notes app, or on blank paper. Record your final thinking here.}

\vspace{4pt}

\renewcommand{\arraystretch}{1}
\begin{tabular}{|>{\bfseries\sffamily\color{tealheader}}p{1.1in}|p{5.7in}|}
\hline
Notice \& Wonder &
\textbf{Model \& Discuss (pg.\ 150b)}

\textbf{Prompt A:} Charlie and Aiko launched a model rocket. It launched from the ground at $t = 0$ seconds and landed at $t = 5$ seconds. Charlie claims the equation is exactly $h(t) = at^2 + bt + c$.

\begin{itemize}[leftmargin=14pt, topsep=2pt, itemsep=0pt]
  \item What do you notice about the launch and landing times?
\end{itemize}

\writespace{30pt}

\begin{itemize}[leftmargin=14pt, topsep=2pt, itemsep=0pt]
  \item What do you wonder about the missing letters $a$, $b$, $c$?
\end{itemize}

\writespace{30pt}

\textbf{Prompt B:} Look at the function $f(x) = x(x - 4)(x + 3)$. Predict the $x$-intercepts using the Zero Product Property.

\writespace{30pt}
\\
\hline
\end{tabular}

\vspace{6pt}

\begin{tabular}{|>{\bfseries\sffamily\color{tealheader}}p{1.1in}|p{5.7in}|}
\hline
Explore &
\textbf{Savvas Example 1 (pg.\ 150)}

\begin{enumerate}[leftmargin=14pt, topsep=2pt, itemsep=2pt]
  \item Use Desmos to graph $f(x) = x(x - 4)(x + 3)$.
  \item Record the exact zeros of the function.
  \item Explain what a zero represents in the rocket situation vs.\ the pure math equation.
  \item Sketch $f(x)$ and pick one ``test point'' in each interval to prove if the graph is above/below the $x$-axis.
\end{enumerate}

\writespace{10pt}
\\
\hline
Discuss &
How is a zero represented in an equation, on a graph, and in a context? Use both functions in your explanation.

\writespace{30pt}
\\
\hline
\end{tabular}

\vspace{6pt}

\begin{cerbox}
In Charlie and Aiko's rocket model, the zeros are $t = 0$ and $t = 5$. Charlie claims he can write the exact equation of the parabola with only that information. Do you agree or disagree? Explain why additional information (like a third point or initial speed) is mathematically required.

\vspace{6pt}
\textbf{Claim:} \blankline

\vspace{4pt}
\textbf{Evidence:} \blankline

\vspace{4pt}
\textbf{Reasoning:} \blankline

\writespace{14pt}
\end{cerbox}

% ===================================================================
%   DAY 2
% ===================================================================
\newpage

\dayheading{2}{Graphing from Factored Form}{Use blank paper for sketches and Desmos for quick checks. Explain every graph using zeros and intervals.}

\vspace{4pt}

\begin{tabular}{|>{\bfseries\sffamily\color{tealheader}}p{1.1in}|p{5.7in}|}
\hline
Compare &
\textbf{Savvas Example 1 (pg.\ 150)}

Compare $f(x) = x(x - 4)(x + 3)$ and $g(x) = -x(x - 4)(x + 3)$.

\begin{itemize}[leftmargin=14pt, topsep=2pt, itemsep=2pt]
  \item What will stay the exact same when graphed?
\end{itemize}

\writespace{24pt}

\begin{itemize}[leftmargin=14pt, topsep=2pt, itemsep=2pt]
  \item What will change? How do you know before graphing?
\end{itemize}

\writespace{24pt}
\\
\hline
\end{tabular}

\vspace{6pt}

\begin{tabular}{|>{\bfseries\sffamily\color{tealheader}}p{1.1in}|p{5.7in}|}
\hline
Build the Graph &
\textbf{Savvas Practice \#12 \& \#13 (pg.\ 156)}

\begin{enumerate}[label=\Alph*., leftmargin=14pt, topsep=2pt, itemsep=2pt]
  \item Use Desmos to check your warm-up prediction.
  \item On blank paper, sketch $f(x) = 3x^3 - 9x^2 - 12x$. \quad\textit{(Hint: factor out the GCF first!)}
  \item Sketch $g(x) = x^3 - 8x^2 + 16x$. \quad(Practice \#13)
\end{enumerate}

\textit{For each function, list the zeros, choose a test point in each interval, and justify the rough shape.}

\writespace{10pt}
\\
\hline
Explain &
Which is more useful when sketching: the zeros or the test points? Defend your answer with at least one example from today.

\writespace{40pt}
\\
\hline
\end{tabular}

\vspace{6pt}

\begin{cerbox}
\textbf{Reason:} If you use zeros to sketch the graph of a polynomial function, how can you verify that your graph is correct? \textit{(Savvas Practice \#6)}

\vspace{6pt}
\textbf{Claim:} \blankline

\vspace{4pt}
\textbf{Evidence:} \blankline

\vspace{4pt}
\textbf{Reasoning:} \blankline

\writespace{14pt}
\end{cerbox}

% ===================================================================
%   DAY 3
% ===================================================================
\newpage

\dayheading{3}{Multiplicity}{Use Desmos to compare graphs and use words like cross, touch, turn, even, odd, and multiplicity.}

\vspace{4pt}

\begin{tabular}{|>{\bfseries\sffamily\color{tealheader}}p{1.1in}|p{5.7in}|}
\hline
Graph \& Compare &
\textbf{Savvas Example 2 (pg.\ 151)}

Graph $y = (x + 2)^2(x - 1)$ and $y = (x + 2)^3(x - 1)$ on Desmos.

\begin{itemize}[leftmargin=14pt, topsep=2pt, itemsep=2pt]
  \item What happens at the zero $x = -2$ in the first graph?
\end{itemize}

\writespace{20pt}

\begin{itemize}[leftmargin=14pt, topsep=2pt, itemsep=2pt]
  \item What happens at $x = -2$ in the second graph?
\end{itemize}

\writespace{20pt}

\begin{itemize}[leftmargin=14pt, topsep=2pt, itemsep=2pt]
  \item What do you predict the exponent (multiplicity) has to do with the behavior?
\end{itemize}

\writespace{20pt}
\\
\hline
\end{tabular}

\vspace{6pt}

\begin{tabular}{|>{\bfseries\sffamily\color{tealheader}}p{1.1in}|p{5.7in}|}
\hline
Create a Rule &
\textbf{Savvas Try It 2b \& Practice \#15}

1.\ For each function below, identify the zero(s), multiplicity, and graph behavior (Crosses or Touches):

\quad a.\ $f(x) = (x^2 + 9)(x - 1)^5(x + 2)^2$

\writespace{20pt}

\quad b.\ $g(x) = x^3 - x^2 - 25x + 25$ \quad\textit{(Hint: factor by grouping!)}

\writespace{20pt}

2.\ Write a rule in words: When multiplicity is odd\ldots\ When multiplicity is even\ldots

\writespace{16pt}

3.\ Test your rule by sketching one new function from your group.
\\
\hline
\end{tabular}

\vspace{6pt}

\begin{generalizebox}
\textbf{Sentence frames:}

When the multiplicity is odd, the graph \blankline\ the $x$-axis.

When the multiplicity is even, the graph \blankline\ at the $x$-axis.
\end{generalizebox}

\vspace{6pt}

\begin{cerbox}
\textbf{Savvas Practice \#7 --- Error Analysis (pg.\ 156)}

\textit{(Hint: Check end behavior AND the number of turning points for a cubic.)}

\vspace{6pt}
\textbf{Claim:} \blankline

\vspace{4pt}
\textbf{Evidence:} \blankline

\vspace{4pt}
\textbf{Reasoning:} \blankline

\writespace{14pt}
\end{cerbox}

% ===================================================================
%   DAY 4
% ===================================================================
\newpage

\dayheading{4}{Real and Complex Zeros}{Finding ALL zeros of a polynomial using graphing, synthetic division, and the quadratic formula.}

\vspace{4pt}

\begin{tabular}{|>{\bfseries\sffamily\color{tealheader}}p{1.1in}|p{5.7in}|}
\hline
Notice \& Predict &
\textbf{Savvas Example 3 (pg.\ 152)}

Graph $f(x) = x^3 - 2x^2 - 2x - 3$.

\begin{itemize}[leftmargin=14pt, topsep=2pt, itemsep=2pt]
  \item The graph shows only one $x$-intercept. What is it?
\end{itemize}

\writespace{16pt}

\begin{itemize}[leftmargin=14pt, topsep=2pt, itemsep=2pt]
  \item Because it's a cubic (degree 3), what do you predict about the other two roots?
\end{itemize}

\writespace{16pt}
\\
\hline
\end{tabular}

\vspace{6pt}

\begin{generalizebox}
\textbf{Three-Step Process}

\textbf{Step 1:} Graph the polynomial to find the real zero.

\textbf{Step 2:} Use Synthetic Division by that real zero to create a smaller quotient.

\textbf{Step 3:} Use the Quadratic Formula $x = \dfrac{-b \pm \sqrt{b^2 - 4ac}}{2a}$ on the quotient to find the remaining complex zeros.
\end{generalizebox}

\vspace{6pt}

\begin{tabular}{|>{\bfseries\sffamily\color{tealheader}}p{1.1in}|p{5.7in}|}
\hline
Guided Example &
\textbf{Find all zeros of $f(x) = x^3 - 2x^2 - 2x - 3$.}

\begin{enumerate}[leftmargin=14pt, topsep=2pt, itemsep=4pt]
  \item Graph it: Real zero at $x =$ \blankline
  \item Synthetic Division: Divide by the real zero. The reduced quotient is $x^2 +$ \blankline $x +$ \blankline.
  \item Quadratic Formula: Solve $x^2 + x + 1 = 0$.
\end{enumerate}

\writespace{30pt}

Final zeros: \rule{3in}{0.4pt}
\\
\hline
Think-Pair-Share &
\textbf{Savvas Try It 3a}

Find all the zeros of $g(x) = x^3 - x^2 - 5x + 5$.

\textit{(Hint: You can use the graph + synthetic division process, OR try factoring by grouping!)}

Student A: find the real zero and set up synthetic division.

Student B: finish the division and set up the Quadratic Formula.

Together: solve for the final roots.

\textit{Language check: Use these words in your explanation --- real, complex, conjugate, quotient.}

\writespace{20pt}
\\
\hline
\end{tabular}

\vspace{6pt}

\begin{cerbox}
\textit{[Refer to Savvas Exit Ticket image]}

\vspace{6pt}
\textbf{Claim:} \blankline

\vspace{4pt}
\textbf{Evidence:} \blankline

\vspace{4pt}
\textbf{Reasoning:} \blankline

\writespace{14pt}
\end{cerbox}

% ===================================================================
%   DAY 5
% ===================================================================
\newpage

\dayheading{5}{Real-World Modeling}{Interpret zeros and sign intervals in context. Use the graph and the equation together.}

\vspace{4pt}

\begin{tabular}{|>{\bfseries\sffamily\color{tealheader}}p{1.1in}|p{5.7in}|}
\hline
Think About the Context &
\textbf{Savvas Example 4 (pg.\ 153)}

\begin{itemize}[leftmargin=14pt, topsep=2pt, itemsep=2pt]
  \item Which zero does \textit{not} make sense in the context of this problem? Why?
\end{itemize}

\writespace{30pt}

\begin{itemize}[leftmargin=14pt, topsep=2pt, itemsep=2pt]
  \item Between what production amounts does Acme make a positive profit?
\end{itemize}

\writespace{30pt}
\\
\hline
\end{tabular}

\vspace{6pt}

\begin{tabular}{|>{\bfseries\sffamily\color{tealheader}}p{1.1in}|p{5.7in}|}
\hline
Waterworks Paddleboards &
\textbf{Savvas Practice \#18 (pg.\ 156)}

Waterworks profit $P(x)$ (in hundreds of dollars) based on $x$ paddleboards sold (in thousands) is modeled by: $P(x) = -3x^3 + 48x^2 - 144x$.

\begin{enumerate}[label=\Alph*., leftmargin=14pt, topsep=2pt, itemsep=2pt]
  \item Factor to find the zeros.
  \item What are the solutions to the equation $P(x) = 0$?
  \item When $P(x)$ is positive, the company earns a profit. Over what interval(s) does the company make a profit?
\end{enumerate}

\writespace{30pt}
\\
\hline
Discuss Domain &
Why must we think about domain in a business problem, even when the algebra allows many more $x$-values?

\writespace{30pt}
\\
\hline
\end{tabular}

\vspace{6pt}

\begin{cerbox}
\textbf{Savvas Practice \#25 --- Make Sense and Persevere}

\textit{[Refer to Savvas textbook image]}

\vspace{4pt}
\textbf{A.}\quad\textbf{Claim:} \blankline\quad\textbf{Evidence:} \blankline\quad\textbf{Reasoning:} \blankline

\writespace{14pt}

\textbf{B.}\quad\textbf{Claim:} \blankline\quad\textbf{Evidence:} \blankline\quad\textbf{Reasoning:} \blankline

\writespace{14pt}

\textbf{C.}\quad\textbf{Claim:} \blankline\quad\textbf{Evidence:} \blankline\quad\textbf{Reasoning:} \blankline

\writespace{14pt}
\end{cerbox}

% ===================================================================
%   DAY 6
% ===================================================================
\newpage

\dayheading{6}{Solving Polynomial Inequalities}{Use zeros and graphs to find where a function is positive or negative.}

\vspace{4pt}

\begin{tabular}{|>{\bfseries\sffamily\color{tealheader}}p{1.1in}|p{5.7in}|}
\hline
Visualizing Inequalities &
\textbf{Savvas Example 6 (pg.\ 154)}

Look at the graph of $f(x) = x^3 - 16x$.

\begin{itemize}[leftmargin=14pt, topsep=2pt, itemsep=2pt]
  \item If $f(x) < 0$, we are looking for where the graph is \textbf{BELOW} the $x$-axis.
  \item Color-code or highlight the exact intervals on the graph where this happens.
\end{itemize}

\writespace{20pt}
\\
\hline
Guided Example &
\textbf{Solve the inequality $x^3 - 16x < 0$ algebraically.}

\begin{enumerate}[leftmargin=14pt, topsep=2pt, itemsep=2pt]
  \item Find the zeros by factoring $x(x^2 - 16) < 0$. Zeros are $0, -4, 4$.
  \item Plot the zeros on a number line and test the intervals: $(-\infty, -4)$, $(-4, 0)$, $(0, 4)$, $(4, \infty)$.
  \item Write the final solution in interval notation.
\end{enumerate}

\writespace{30pt}
\\
\hline
Partner Practice &
\textbf{Savvas Try It 6a \& 6b}

\begin{enumerate}[label=\Alph*., leftmargin=14pt, topsep=2pt, itemsep=2pt]
  \item Solve: $2x^3 + 12x^2 + 12x < 0$
  \item Solve: $(x^2 - 1)(x^2 - x - 6) > 0$
\end{enumerate}

\textit{Use blank paper for your number line and sign testing. Check your answers by graphing on Desmos.}

\writespace{20pt}
\\
\hline
\end{tabular}

\vspace{6pt}

\begin{generalizebox}
\textbf{Explanation frame:}

The solution includes the intervals \blankline\ and \blankline\ because the graph is \blankline\ the $x$-axis in those regions. We use brackets instead of parentheses because the inequality is \blankline.
\end{generalizebox}

\vspace{6pt}

\begin{cerbox}
A student solves $2x^3 + 12x^2 + 12x > 0$ and writes the answer as $[-6, -1] \cup [0, \infty)$. Evaluate their answer. Is their interval notation correct? Explain the difference between using parentheses $(\;)$ and brackets $[\;]$ in inequalities.

\vspace{6pt}
\textbf{Claim:} \blankline

\vspace{4pt}
\textbf{Evidence:} \blankline

\vspace{4pt}
\textbf{Reasoning:} \blankline

\writespace{14pt}
\end{cerbox}

% ===================================================================
%   DAY 7
% ===================================================================
\newpage

\dayheading{7}{Synthesis and Critique}{Use this page to argue, critique, and connect the big ideas from the whole unit.}

\vspace{4pt}

\begin{tabular}{|>{\bfseries\sffamily\color{tealheader}}p{1.1in}|p{5.7in}|}
\hline
SAT/ACT Practice &
\textbf{Savvas Practice \#29 (pg.\ 157)}

Without using a graphing calculator, determine which of the 4 graphs provided represents $f(x) = x^3 + x^2 - 4x$.

\textit{Explain to your group how end-behavior and factoring helped you eliminate the wrong choices.}

\writespace{40pt}
\\
\hline
Performance Task &
\textbf{Savvas Practice \#30: Venetta's Deli (pg.\ 157)}

Venetta opened several deli sandwich franchises in the year 2000. Her profit $P(t)$, in hundreds of dollars, after $t$ years can be modeled by $P(t) = t^3 + t^2 - 6t$.

\textbf{Part A:} Sketch a graph of the function.

\hfill\gridpaper{5}{3.5}\hfill\mbox{}

\textbf{Part B:} Based on the model, during what years did Venetta NOT make a profit?

\writespace{26pt}

\textbf{Part C:} If the model is appropriate, predict the amount of profit Venetta will receive from her franchises in the year 2020.

\writespace{26pt}
\\
\hline
Listen \& Respond &
Listen to another group's explanation. Add one question, one counterexample, or one way to strengthen the evidence.

\writespace{24pt}
\\
\hline
\end{tabular}

\vspace{6pt}

\begin{cerbox}
\textbf{Critique the Model}

Based on your answer to Part C of the Venetta problem, you predicted a specific profit for the year 2020. Critique the limits of this mathematical model. Is it realistic to assume a cubic function will perfectly model business profit forever? Why or why not?

\vspace{6pt}
\textbf{Claim:} \blankline

\vspace{4pt}
\textbf{Evidence:} \blankline

\vspace{4pt}
\textbf{Reasoning:} \blankline

\writespace{14pt}
\end{cerbox}

% ===================================================================
%   DAY 8
% ===================================================================
\newpage

\dayheading{8}{Mastery Check and Reflection}{Use the assessment to show what you can do, then choose the station that best helps you keep growing.}

\vspace{4pt}

\begin{tabular}{|>{\bfseries\sffamily\color{tealheader}}p{1.1in}|p{5.7in}|}
\hline
Error Analysis &
\textbf{Savvas Practice \#27 (pg.\ 157)}

The height of a rectangular storage box is less than both its length and width. The function $f(x) = x^3 + 2x^2 - 3x$ represents the volume, where $x$ is the width in feet.

\begin{enumerate}[label=(\alph*), leftmargin=14pt, topsep=2pt, itemsep=2pt]
  \item Find the factored form.
  \item Find the zeros.
  \item What do the other two factors represent?
  \item Find dimensions when volume is 10 ft$^3$.
\end{enumerate}

\writespace{30pt}
\\
\hline
Mastery Check &
\textbf{Savvas Lesson Quiz (pg.\ 157)}

\textit{Do your best on these 3 problems without graphing technology.}

\writespace{120pt}
\\
\hline
\end{tabular}

\vspace{6pt}

\begin{tabular}{|>{\bfseries\sffamily\color{tealheader}}p{1.1in}|p{5.7in}|}
\hline
Choose a Station &
\textbf{Savvas Differentiated Resources}

\textbf{Station 1 (Reteach):} Work with the teacher on ``Reteach to Build Understanding'' (focus on factoring and sign charts).

\vspace{4pt}
\textbf{Station 2 (Mathematical Literacy):} Work on the Vocabulary worksheet matching graphs to roots, multiplicities, and terms.

\vspace{4pt}
\textbf{Station 3 (Enrichment):} Complete the advanced problem-solving task predicting higher-degree polynomial behaviors.
\\
\hline
\end{tabular}

\vspace{6pt}

\begin{cerbox}
Reflect on your learning. When trying to graph or solve a polynomial, which representation do you find most helpful and why: Factoring/Algebra, Desmos/Visual Graphing, or Synthetic Division?

\vspace{6pt}
\textbf{Claim:} \blankline

\vspace{4pt}
\textbf{Evidence:} \blankline

\vspace{4pt}
\textbf{Reasoning:} \blankline

\writespace{20pt}
\end{cerbox}

\end{document}
